% !TEX encoding = UTF-8
% !TEX TS-program = pdflatex
% !TEX root = ../tesi.tex

% Sommario
\cleardoublepage
\phantomsection
\pdfbookmark{Sommario}{Sommario}
\begingroup
\let\clearpage\relax
\let\cleardoublepage\relax
\let\cleardoublepage\relax

\chapter*{Sommario}

\noindent Il presente documento descrive il lavoro svolto durante il periodo di stage, effettuato presso l'azienda Moku S.r.l. tra Maggio e Luglio 2022. Lo scopo di tale attività era progettare e sviluppare un'applicazione \textit{mobile} multipiattaforma per la gestione di tour cicloturistici. Gli obiettivi da raggiungere durante il tirocinio erano molteplici. \\

\noindent In primo luogo era richiesto di sviluppare la base dell'applicazione \textit{mobile}, seguendo gli stili e i \gls{mockupg} forniti dal cliente, in secondo luogo veniva richiesto di implementare la funzionalità di condivisione delle posizioni tra i partecipanti di un tour e la visualizzazione in dettaglio di un tour e delle sue caratteristiche, infine era richiesta l'implementazione delle schermate di autenticazione. \\

\noindent Il lavoro si è concentrato principalmente sulle precedenti attività, riuscendo anche a completare schermate opzionali e non richieste. Alla conclusione dello stage, tutte le funzionalità previste per il tirocinio erano concluse. \\

\noindent Il documento è suddiviso in quattro capitoli:
\begin{itemize}
    \item \hyperref[cap:cap1]{Il primo capitolo}: in questo capitolo viene introdotto il contesto aziendale del tirocinio, fornendo una breve paronamica sull'azienda e la sua storia. Vengono infine introdotti il progetto proposto, le sue caratteristiche e gli obiettivi previsti per tutto lo svolgimento del tirocinio;
    \item \hyperref[cap:cap2]{Il secondo capitolo}: in questo capitolo vengolo illustrati i processi di sviluppo e gli strumenti adottati per affrontare il progetto;
    \item \hyperref[cap:cap3]{Il terzo capitolo}: in questo capitolo ci si addentra nello svolgimento del progetto, affrontando le varie fasi percorse, quali l'analisi dei requisiti, la progettazione, la codifica e la validazione del prodotto;
    \item \hyperref[cap:cap4]{Il quarto capitolo}: in questo capitolo viene attuata una valutazione retrospettiva del tirocinio svolto, andando a stilare gli obiettivi raggiunti e le conoscenze acquisite. Viene infine ritagliato uno spazio per le considerazioni personali sull'esperienza del tirocinio.
\end{itemize}

\paragraph{Convenzioni tipografiche}
\begin{itemize}
	\item gli acronimi, le abbreviazioni e i termini ambigui o di uso non comune menzionati vengono definiti nel glossario, situato alla fine del presente documento;
	\item i termini in lingua straniera o facenti parti del gergo tecnico sono evidenziati con il carattere \textit{corsivo}.
\end{itemize}

\endgroup
\vfill

